% !TeX root = ../papers/nc-p3.tex
\section{非半単純な相補逆行列の全補完と閉軌道}
\label{sec:split-jordan-transport}
前節の二次形式に関する行列の同一視を使い、
$A=M^\dagger M$, $S=\operatorname{tr}A$, $d=\det M$ と置く。
固有値を $x_i^2$ と書き、$\prod x_i=d$ とする。
\begin{theorem}\label{thm:split-jordan-transport}
$\det M\ne0$ とする。可逆二平面分解の窓に $M$ が現れる必要十分条件は、
$x$ が $f_0f_1\ne0$ を満たす前節の曲面・直線・六次式のいずれかの付け替えであることである。
重固有値ごとのJordan分割への追加条件はない。非常に偶な場合（一つ一つの鎖長が全て偶数）
の直交基底の向きの双方でも成立する。全ての補完の分解付き同型は
定理\ref{thm:split-inverse-transporter}の共役商である。
\end{theorem}
\begin{proof}[二つの逆テンソルの単一行列による表示]
標準的な二側の形式を互いに同一視する。$f_0=d$ として
\[
 16 f_1=S^2-2\operatorname{tr}(A^2)+8d,\qquad
 F(z)=Sz-2z^2+2d,\qquad G(z)=F(z)^2/z .
\]
実際、最初の等式はBell基底で前節の四一次因子を展開した等式であり、
双方が特殊直交群二つの不変式なので、単純スペクトルの稠密性により全行列に成立する。
勾配を形式で同一視すると、$\nabla f_0=d M A^{-1}$ と
$4\nabla f_1=M A^{-1}F(A)$ を得る。従って二つの逆テンソルは射影的に
$M_0=M A^{-1},\ M_1=M A^{-1}F(A)$ と書ける。その右Gram行列は
$A^{-1},G(A)$ である。第二式の勾配と二つの切断の逆転置との対応は
\path{computations/verify_split_jordan_transport.py} でも通常基底で確認する。
$\tau$ は行側と列側を交換する。転置後も可逆行列の左右Gramは相似なので、
移送があれば平方固有値の多重集合・行列式は半単純化した二つの逆テンソルでも合う。
これが半単純な $P$ の局所同値の必要十分条件だったので、前節のスペクトル条件は必要である。

自己随伴行列は各一般固有空間の中で、鎖のGram行列が反対角単位行列となる
Jordan鎖の直交和にできる。実際、最大鎖長を $m$ とし、冪零部 $N$ について
$\langle v,N^{m-1}v\rangle\ne0$ と取れば、その巡回空間は非退化である。
$k[N]/N^m$ の単元の平方根で生成元を変更して最後のGram係数だけを1にし、
直交補に帰納する。従って同じJordan型の自己随伴行列は直交共役である。
一つでも奇数長の鎖があればその上だけの符号で共役の行列式を変更できる。

各許容スペクトル上で $M=O L$ と取る。$L$ は上の各鎖で $A$ の単元平方根を取った
可逆自己随伴行列である。奇数長の鎖があれば符号を選んで $\det O=1$ にできる。
偶数長だけなら $\det L=\prod_\nu\nu^{\dim E_\nu/2}$ である（$E_\nu$ は一般固有空間）。
直線上のスペクトルは $(1,1,r,r)$ に $d=r$、曲面上で全多重度が偶数となるのは
これと重なる $(1,1,-1,-1)$ に $d=-1$ である。従ってここでも $O$ は特殊直交である。
座標・共通倍率の変更はこの存在性を変えない。
$M_1=O L A^{-1}F(A)$ の右因子は自己随伴なので、その随伴も二側の特殊直交作用で
$M_1$ に戻る。よって $\tau$ に伴う転置を許した相補逆行列の比較を、
$M_0,M_1$ そのものの二側特殊直交同値として判定できる。

直線上では $\nu\in\{1,r\}$ で $F(\nu)=4r$, $G'(\nu)=-16r\ne0$ である。
従って $G(A)$ は $(4r)^2 A^{-1}$ と同じ固有値・Jordan分割を持つ。
曲面では $x=(1,-uv,u,v)$、$r=u^2,s=v^2,C=(r-1)(s-1)\ne0$、
$d=-rs$ とおく。$\nu=(1,rs,r,s)$ の順に
$F(\nu)/\nu=(-C,-C,C,C)$。$G(\nu)=C^2\nu$ は、$A^{-1}$ の固有値を
$\nu\mapsto rs/\nu$ で対応させた後 $rsC^2$ 倍したものである。
$1$ と $rs$ または $r$ と $s$ が重なる場合だけ非自明な鎖が可能で、
どちらもこの置換がその重固有値を固定する。
$r=s$ の $\nu=r$ で $G'(r)=(r-1)^4$、$rs=1$ の $\nu=1$ で
$G'(1)=-C(r+r^{-1}-2)$ なので鎖の長さは保たれる。
単純固有値の六次式の場合は既に前節で解いた。
上で得たGram倍率を $c^2$ とすれば、直線と曲面の双方で $\det F(A)=c^4$ である。

こうして $M_1$ と $cM_0$ の行列式が等しく、右Gramが直交共役となった。
必要なのは特殊直交共役である。奇数鎖の場合は先の符号を使う。
偶数鎖のみなら、長さ4が一つのときは冪零部の平方の核、
二つの長さ2のときは各一般固有空間内の冪零部の核の直和を取る。
これは双方で同じ極大全等方二平面である。導関数が零でないので同じ核を持ち、
スペクトルの並べ替えもこの合計を変えない。この平面を保つ直交行列の対角ブロックは
$B,B^{-\mathsf t}$ なので行列式は1である。
右の同値行列が特殊直交と分かれば、左の行列もGramの等式と行列式の一致から特殊直交である。
\end{proof}

\begin{theorem}\label{thm:split-polystable-good-limits}
可逆二平面分解の全窓は半安定で、唯一の閉軌道も同じ可逆二平面分解の
標準Koszul左右AS正則窓である。非半単純な $P$ のどの補完も、同じ平方固有値と
行列式を保つ半単純な $P$ の補完を軌道閉包内に持つ。
\end{theorem}
\begin{proof}
標準化した右Jordan鎖に $1-m,3-m,\ldots,m-1$ の重みをつけ、
$N$ が重みを上げる向きとする。反対角Gramよりこれは直交一径数である。
誘導される減少フィルトレーションは鎖の選択に依存しない。実際、中心0の冪零重み
フィルトレーションは、各Jordan鎖で上記の重みになるものとして、核と像の交叉の和で書ける。
具体的に重み $\ge j$ の部分は $\sum_{a\ge0}\ker N^{a+1}\cap
\operatorname{im}N^{\max(0,a+j)}$ である。一本の長さ $m$ の鎖で、先頭から
$h$ 番目（番号0始まり）のベクトルについての条件が $2h-(m-1)\ge j$ になる。
一般固有空間ごとにこれを取り、左側には $O$ により同じ一径数を運ぶ。
$L$ はフィルトレーションを両向きに保つので $M$ も左右のフィルトレーションを対応させる。
次数付きの $M$ は可逆で右Gramの冪零部は消える。

$M_0,M_1$ の各Gram上でも、非自明鎖での導関数の非消滅から冪零部の全ての核・像は
対応する元のものと同じである。従って左右のこのフィルトレーションは両逆テンソルから
同じ規則で再現できる。移送 $v\in\tau D$ はこれらを行と列の交換を伴って保つ。
上の一径数を $\lambda(t)\in D$ とすれば $\lambda(t)v\lambda(t)^{-1}$ は群内に極限を持つ。
実際、フィルトレーションを保つ行列は重みを下げる成分がなく、逆にも同じことが成立する。
二側の特殊直交一径数を射影二文字因子に戻し、必要なら $t$ の有限冪を取って線形に持ち上げる。
第一の四対応を固定した前節の構成で、交代成分と第二の対応がともにこの次数付きに特殊化する。
各局所射影行列は可逆な極限を持ち、特殊線形持上げでは発散する共通倍率はない。
$f_0,f_1$ も非零のままであり、得られた窓は半単純な $P$ の補完である。

さらに、全基底群内で分解を保つ中心化群でのアフィン軌道閉包を考える。
六つの二重次数の成分ごとに、可逆な二対二切断の四次行列式を一つ取る。
高位位置集合 $I$ に対応する行列式は四位置の使用した二平面の行列式の二乗の積で変換する。
従って $I$ と補集合の二行列式の積は、各位置の行列式を1とした中心化群に不変である。
両交代成分では両切断を使ってよい。これらの不変式は元の窓で非零なので、
閉包内でも各境界成分は非退化二形式二つの積のままであり、両交代成分の必要な切断も可逆である。
切断階数の上界は極限へ延び、可逆な交差ブロックで下界も保たれる。
従ってこの閉包内の中心化群の閉軌道は零でなく、同じ正則枝に属する。
その閉軌道でHermiteノルムが最小になるように基底を取れば、各位置の
二半分はトーラスの異なる重みでGram行列の交差ブロックが零であり、
残る全ての跡零方向は中心化群の極小条件に含まれる。
よって全基底群でも閉軌道であり、還元と最小ノルムの議論は標数零の代数閉体に戻る。
\end{proof}
